\documentclass[11pt]{article}
\usepackage{amsmath,amssymb}
\usepackage{hyperref}
\usepackage{geometry}
\usepackage{booktabs}
\geometry{margin=1in}

\title{The Silence-First Framework (SSF) v0.2:\\An Experimental Pre-Inference Admission Regulator}
\author{Muze-X \\ \small{Research direction: Muse-IRS}}
\date{Public preprint release: 2026-09-17}

\begin{document}
\maketitle

\begin{abstract}
We present a reconstruction of the Silence-First Framework (SSF), an
experimental pre-inference admission regulator. The central idea is narrow:
non-action should be representable as a valid computational state.

The initial repository mixed implementation, physical interpretation,
universal-invariant language, and benchmark claims. Version 0.2 separates
these levels. It defines an abstract recoverable state, a filtered state, a
normalized load ratio, and a deterministic three-state gate returning
\texttt{ACTIVE}, \texttt{SILENCE}, or \texttt{DEFER}. We do not claim a
universal causal invariant, production readiness, or general performance
improvement. Those questions are left to explicit adapters and reproducible
experiments.
\end{abstract}

\section{Motivation}

A common software path is
\[
\text{trigger} \rightarrow \text{execution}.
\]
SSF investigates a different architecture:
\[
\text{trigger} \rightarrow \text{pre-action regulator}
\rightarrow \{\text{ACTIVE},\text{SILENCE},\text{DEFER}\}.
\]
The regulator is intentionally independent of semantic content moderation,
ethical reasoning, and model-output alignment.

\section{Epistemic Scope}

The reference implementation is a model, not a physical law. We distinguish:
\begin{itemize}
\item implementation relations defined by the code;
\item modeling hypotheses;
\item predictions that require experimental protocols;
\item results obtained under explicit tests;
\item unknown relations that remain open.
\end{itemize}

A test that confirms implementation behavior does not validate a physical or
system-wide interpretation.

\section{State Model}

Let $\Delta\phi_t$ denote an adapter-supplied abstract increment. The v0.2
kernel uses
\[
\phi_{t+1}=\lambda\phi_t+\Delta\phi_t,
\]
where $\lambda\in[0,1]$ is a retention parameter, followed by
\[
\phi^c_{t+1}=\alpha\phi_{t+1}+(1-\alpha)\phi^c_t,
\]
with $\alpha\in(0,1]$.

Given a positive configured capacity $\Omega$, we define the model load ratio
\[
q_t=\frac{|\phi^c_t|}{\Omega}.
\]
This quantity is dimensionless only when the adapter and capacity definition
make that normalization meaningful. It is not presented as a universal
invariant.

\section{Admission Policy}

The current reference policy uses an activity threshold $\varepsilon$, an
overload threshold $q_{\max}$, a recovery threshold $q_{\mathrm{rec}}$, and
the state variables above, with
\[
0\le q_{\mathrm{rec}}<q_{\max}.
\]

The decision logic is:
\begin{enumerate}
\item if the regulator is already overloaded and $q_t>q_{\mathrm{rec}}$,
return \texttt{DEFER};
\item if $q_t\ge q_{\max}$, enter overload and return \texttt{DEFER};
\item if no previous filtered sample exists, return \texttt{ACTIVE}; the local
change is unknown and is not treated as evidence for silence;
\item once a baseline exists, if
$|\phi^c_t-\phi^c_{t-1}|\le\varepsilon$, return \texttt{SILENCE};
\item otherwise return \texttt{ACTIVE}.
\end{enumerate}

The separate entry and recovery thresholds provide hysteresis. This is an
engineering hypothesis to test, not an asserted optimum.

\section{Interpretation of States}

\texttt{ACTIVE} means that execution is admitted by the current policy. It
does not imply semantic correctness or safety.

\texttt{SILENCE} means that execution is not activated because local filtered
change is below the configured threshold after a baseline exists. It does not
imply that an input has no semantic importance.

\texttt{DEFER} means that execution is temporarily inhibited while the model
state remains above the configured recovery region. It is distinct from a
permanent refusal.

\section{Revision of the Historical Invariant Claim}

The initial repository defined
\[
L=\frac{\Phi_c\pi}{\Omega}
\]
and described $L$ as a universal causal load invariant spanning informational,
thermodynamic, stochastic, and temporal regimes.

The current repository does not demonstrate a transformation class under
which this quantity is invariant, nor a general procedure that makes those
heterogeneous domains commensurable. Version 0.2 therefore removes the
universal-invariant status from the operational core. The historical formula
is retained only as part of the project genealogy and has status
\emph{not established}.

\section{Experimental Program}

A domain-level evaluation must specify:
\begin{enumerate}
\item a measurable workload or system quantity;
\item an adapter mapping that quantity to $\Delta\phi_t$;
\item units, normalization, sampling interval, and valid range;
\item a baseline admission policy;
\item SSF configuration;
\item an oracle or evaluation rule for acceptable activity and silence;
\item reproducible traces and metrics.
\end{enumerate}

Candidate metrics include execution count, resource use, latency, false
silence, false activity, duration of deferral, and recovery time.

The earlier statement that SSF reduced inference calls by approximately
35--65\% without semantic degradation is not treated as an established result
in v0.2 because the current repository does not contain the full reproducible
benchmark needed to support that claim.

\section{Limitations}

The present kernel does not establish:
\begin{itemize}
\item universal causal meaning of $\phi$ or $q$;
\item physical, thermodynamic, or quantum interpretation;
\item optimal thresholds;
\item semantic quality preservation;
\item safety-critical suitability;
\item production-level performance gains.
\end{itemize}

\section{Conclusion}

SSF v0.2 reduces the framework to a falsifiable engineering proposition:
\begin{quote}
A system may need the capacity not to act now in order to preserve the
capacity to act later.
\end{quote}
Whether this proposition is useful in a given domain must be established by
measurable adapters, explicit baselines, and reproducible tests.

\section*{Collaborative Contribution}

Human research direction and publication accountability: Muze-X / Muse-IRS.
ChatGPT (OpenAI) contributed through AI-assisted formalization, scientific
reconstruction, counter-analysis, falsification design, documentation, and
manuscript development. This disclosure records the collaborative construction
process and is not independent scientific certification.

\bibliographystyle{plain}
\bibliography{references}
\end{document}
